\documentclass[11pt]{amsart}
\pagestyle{plain}
\raggedbottom
\usepackage[dvipsnames]{xcolor}
\usepackage[utf8]{inputenc}
\usepackage{amssymb,amsmath,amsthm,amsfonts,fullpage,float,latexsym,bbm,microtype,cite,fancyvrb}
\usepackage{tikz}
\usepackage{enumerate}
\usetikzlibrary{decorations.pathreplacing}
\usepackage{hyperref}
\hypersetup{colorlinks=true, linkcolor=blue, citecolor=magenta, filecolor=magenta, urlcolor=magenta}
\usepackage{bm,mathrsfs}





\makeatletter
\newtheorem*{rep@theorem}{\rep@title}\newcommand{\newreptheorem}[2]{%
\newenvironment{rep#1}[1]{%
\def\rep@title{\bf #2 \ref{##1}}%
\begin{rep@theorem}}%
{\end{rep@theorem}}}
\makeatother
\newreptheorem{theorem}{Theorem}

\newtheorem{theorem}{Theorem}
\newtheorem{proposition}[theorem]{Proposition}
\newtheorem{claim}[theorem]{Claim}
\newtheorem{conjecture}[theorem]{Conjecture}
\newtheorem{lemma}[theorem]{Lemma}
\newtheorem{corollary}[theorem]{Corollary}
\theoremstyle{definition}
\newtheorem{remark}[theorem]{Remark}
\newtheorem{definition}[theorem]{Definition}
\newtheorem{observation}[theorem]{Observation}
\newtheorem{example}[theorem]{Example}
\newtheorem{problem}[theorem]{Problem}
\newtheorem{comment}[theorem]{Comment}


\renewcommand{\emptyset}{\varnothing}
\renewcommand{\sl}{\mathfrak{sl}}
\newcommand{\0}{\emptyset}
\newcommand{\Q}{\mathbb{Q}}
\newcommand{\R}{\mathbb{R}}
\newcommand{\Z}{\mathbb{Z}}
\newcommand{\C}{\mathbb{C}}
\newcommand{\N}{\mathbb{N}}

\DeclareMathOperator{\Hilb}{Hilb}
\DeclareMathOperator{\Frob}{Frob}
\DeclareMathOperator{\stair}{stair}
\DeclareMathOperator{\Stir}{Stir}
\DeclareMathOperator{\SW}{SW}
\DeclareMathOperator{\SF}{SF}
\DeclareMathOperator{\sminv}{sminv}
\DeclareMathOperator{\Split}{Split}
\DeclareMathOperator{\st}{st}
\DeclareMathOperator{\wt}{wt}
\DeclareMathOperator{\rank}{rank}
\DeclareMathOperator{\Char}{Char}
\DeclareMathOperator{\GL}{GL}
\DeclareMathOperator{\IrrChar}{IrrChar}
\DeclareMathOperator{\Cat}{Cat}
\DeclareMathOperator{\rev}{rev}
\DeclareMathOperator{\Sym}{Sym}
\DeclareMathOperator{\SST}{SST}
\DeclareMathOperator{\diag}{diag}
\DeclareMathOperator{\Res}{Res}
\DeclareMathOperator{\ch}{char}
\DeclareMathOperator{\OSP}{OSP}
\DeclareMathOperator{\inv}{inv}
\DeclareMathOperator{\sgn}{sgn}
\DeclareMathOperator{\Inv}{Inv}
\DeclareMathOperator{\blocks}{blocks}


\DeclareRobustCommand{\qbinom}{\genfrac[]{0pt}{}}





\begin{document}

\title{1st proof problem}

\author[S. Corteel]{Sylvie Corteel}
\address{Department of Mathematics\\
         University of California, Berkeley, CA, USA}
\email{corteel@berkeley.edu}

\author[J. Lentfer]{John Lentfer}
\address{Department of Mathematics\\
         University of California, Berkeley, CA, USA}
\email{jlentfer@berkeley.edu}


\maketitle

\section{The Problem}

Let $m$ and $n$ be positive integers and let $R_n^{(m)}$ be the coinvariant algebra with $m$ sets of $n$ variables.

The Hilbert series of an $m$-multigraded algebra $A$ is
$$
{\rm Hilb}(A;q_1,\ldots ,q_m)= \sum_{i_1,\ldots,i_m \geq 0} \dim(A_{(i_1,\ldots,i_m)})q_1^{i_1} \cdots q_m^{i_m},
$$
where $A_{(i_1,\ldots,i_m)}$ denotes a multigraded component of the algebra, with grading variables $q_1,\ldots, q_m$.

Let $s_\lambda(q_1,\ldots ,q_m)$ denote a Schur polynomial. For $R_n^{(m)}$, it is known that its Hilbert series may be expressed by 
$$
{\rm Hilb}(R_n^{(m)};q_1,\ldots ,q_m)= \sum_{\lambda}
c_{\lambda}(n)s_\lambda(q_1,\ldots ,q_m),
$$
for some nonnegative universal series coefficients $c_\lambda(n)$ which do not depend on $m$.

For any $n$, give a combinatorial interpretation of the coefficients $c_{\lambda}(n)$ when $\lambda$ is a hook shape partition.

\section{The Solution}
\subsection{Background}

The \textbf{coinvariant ring with $m$ set of variables} is defined by
\begin{equation}\label{eq:m-coinv} R_n^{(m)} = \C[x_1^{(1)},\ldots,x_n^{(1)},\ldots, x_1^{(m)}, \ldots,x_n^{(m)}]/\langle \C[x_1^{(1)},\ldots,x_n^{(1)},\ldots, x_1^{(m)}, \ldots,x_n^{(m)}]_+^{\mathfrak{S}_n} \rangle,\end{equation}
where $\C[x_1^{(1)},\ldots,x_n^{(1)},\ldots, x_1^{(m)}, \ldots,x_n^{(m)}]_+^{\mathfrak{S}_n}$ denotes the polynomials without constant term, which are invariant under the diagonal action of $\mathfrak{S}_n$ permuting each set of variables.

The multigraded Hilbert series of $R_n^{(m)}$ is defined by 
\begin{equation} \Hilb(R_n^{(m)}; q_1,\ldots,q_m) := \sum_{r_1,\ldots,r_m \geq 0}\dim \left((R_n^{(m)})_{r_1,\ldots,r_m} \right)q_1^{r_1}\cdots q_m^{r_m}.\end{equation}
Following from the work of Bergeron \cite{Bergeron2020}, this Hilbert series may be written in the form
\begin{equation} \Hilb(R_n^{(m)}; q_1,\ldots,q_m) = \sum_{\lambda, \ \ell(\lambda) \leq m} c_\lambda(n) s_\lambda(q_1,\ldots,q_m),
\end{equation}
for some nonnegative integer coefficients $c_\lambda(n)$, which are independent of $m$. 
It is important to note that while they are independent of $m$, $s_\lambda(q_1,\ldots,q_m)$ can evaluate to $0$ if $\lambda$ has too many parts relative to $m$.
Given a fixed $n \geq 1$, all of the possible $c_\lambda(n)$ coefficients that could possibly appear do appear once $m \geq n$, by \cite{Bergeron2020}. Thus we can reduce to the case where $m=n$:
\begin{equation} \Hilb(R_n^{(n)}; q_1,\ldots,q_n) = \sum_{\lambda, \ \ell(\lambda) \leq n} c_\lambda(n) s_\lambda(q_1,\ldots,q_n).
\end{equation}

Now in the problem, we are asked to give an explicit combinatorial formula for the $c_\lambda(n)$ when $\lambda$ is a hook shape partition. 
A partition $\lambda$ which satisfies $\lambda_2 \leq 1$ is called a \textbf{hook shape}. 
By diagonal supersymmetry (see Proposition~\ref{prop:ds}), conjectured by \cite{Bergeron2020} and proven by \cite{Lentfer-Supersymmetry}, we are able to convert this to a problem about the superspace coinvariant ring. This is because all of the hook shapes $\lambda$ which satisfy $\ell(\lambda) \leq n$ can be determined from this setting.

The \textbf{superspace coinvariant ring} is defined by
\begin{equation}\label{eq:super-coinv} R_n^{(1,1)} = \C[x_1,\ldots,x_n, \theta_1, \ldots,\theta_n]/\langle \C[x_1,\ldots,x_n, \theta_1, \ldots,\theta_n]_+^{\mathfrak{S}_n} \rangle,\end{equation}
where $\C[x_1,\ldots,x_n, \theta_1, \ldots,\theta_n]_+^{\mathfrak{S}_n}$ denotes the polynomials without constant term, which are invariant under the diagonal action of $\mathfrak{S}_n$ permuting each set of variables.
The $x_i$ are commutative variables, and the $\theta_i$ are anticommutative variables.
Commuting variables commute with all variables, while anticommuting variables anticommute with all anticommuting variables, that is, $\theta_i\theta_j = -\theta_j \theta_i$.

The superspace coinvariant ring $R_n^{(1,1)}$ is bigraded, where each $x_i$ contributes to the $q$-degree and each $\theta_i$ contributes to the $u$-degree. 
Let $(R_n^{(1,1)})_{r,s}$ denote the subspace of $R_n^{(1,1)}$ which has $q$-degree $r$ and $u$-degree $s$.
Then the bigraded Hilbert series is defined by 
\begin{equation} \Hilb(R_n^{(1,1)}; q; u) := \sum_{r,s \geq 0}\dim \left((R_n^{(1,1)})_{r,s} \right)q^{r} u^{s}.\end{equation}

Let $[k]_q$ denote the $q$-integer $1+q+\cdots+q^{k-1}$. Define the $q$-factorial number $[k]_q!$ by $[k]_q [k-1]_q \cdots [2]_q [1]_q$, and let $[0]_q! := 1$. Then define the $q$-binomial coefficients by 
\begin{equation}
    \qbinom{n}{k}_q = \frac{[n]_q!}{[k]_q! [n-k]_q!}.
\end{equation}

Ordered set partitions of $\{1,\ldots,n\}$ with $k$ blocks are sequences of nonempty sets $w = (S_1/ S_2/ \cdots / S_k)$ such that the disjoint union of the sets $S_i$ is $\{1,\ldots,n\}$. 
We denote the set of all ordered set partitions of $\{1,\ldots,n\}$ into $k$ blocks by $\OSP(n,k)$.
Following Sagan--Swanson \cite[Section 5.1]{SaganSwanson2024}, define the inversion set on ordered set partitions by
\begin{equation}
    \Inv(w) = \{(s,S_j)\, | \, s \in S_i \text{ for some } i < j \text{ and } s \geq \min S_j\},
\end{equation}
and the inversion statistic by $\inv(w) = \# \Inv(w)$.

Recall the \textbf{$q$-Stirling numbers of the second kind} $S[n,k]$, which are defined by the initial conditions $S[0,k] = \delta_{0,k}$ and $S[n,0] = \delta_{n,0}$, along with the recurrence relation
\begin{equation}\label{eq:stirling}
    S[n,k] = S[n-1,k-1] + [k]_q S[n-1,k]
\end{equation}
for $0 \leq k \leq n$.

Rhoades and Wilson proved the following explicit formula for the bigraded Hilbert series of $R_n^{(1,1)}$.
\begin{theorem}[\!\!\cite{RhoadesWilson2023}]\label{thm:hilbert}
    For any $n \geq 1$,
    \begin{equation}\label{eq:hilbert-rhoades-wilson}
        \Hilb(R_n^{(1,1)};q;u) = \sum_{k=0}^n [k]_q!S[n,k] u^{n-k}.
    \end{equation}
\end{theorem}


Denote a Schur function by $s_\lambda(q_1,\ldots,q_k)$ and a skew Schur function by $s_{\lambda/\mu}(q_1,\ldots,q_k)$.
A \textbf{super Schur function} $s_\lambda(q_1,\ldots,q_k/u_1,\ldots,u_j)$ is defined by 
\begin{equation}
    s_\lambda(q_1,\ldots,q_k/u_1,\ldots,u_j) = \sum_{\nu \subseteq \lambda} s_\nu(q_1,\ldots,q_k)s_{\lambda'/\nu'}(u_1,\ldots,u_j),
\end{equation}
where $\lambda'$ denotes the transpose of the partition $\lambda$ (see for example \cite[Sections A.2.2 and 3.2]{ChengWangBook}).

We recall diagonal supersymmetry for Hilbert series of coinvariant rings, specialized to the superspace coinvariant ring.
Let $P(k,j,n)$ denote the set of all partitions which satisfy $\lambda_{k+1} \leq j$ and the length $\ell(\lambda) \leq n$.

\begin{proposition}[\!\!\cite{Lentfer-Supersymmetry}]\label{prop:ds}
    Fix a positive integer $n$. For partitions $\lambda$ with $\ell(\lambda) \leq n$, there exist nonnegative integer coefficients $c_\lambda(n)$ such that the bigraded Hilbert series of $R_n^{(1,1)}$ is
\begin{equation}\label{eq:coeffs-Hilbert}
    \Hilb(R_n^{(1,1)}; q;u) = \sum_{\lambda \in P(1,1,n)} c_{\lambda}(n) s_\lambda(q/u).
\end{equation}
\end{proposition}
Note that $P(1,1,n)$ consists of all hook shape partitions of length at most $n$. This is a key insight in the simplification of the problem.


The super Schur function has a simple formula for hook shape partitions.

\begin{lemma}\label{lem:schur-1-1}
    If $\lambda$ is $(a,1^b)$ for some $a \geq 1$ and $b \geq 0$, then $s_\lambda(q/u) = q^au^b + q^{a-1}u^{b+1}$. 
    If $\lambda$ is $\varnothing$, then $s_\lambda(q/u) = 1$.
    If $\lambda$ is not contained in a hook shape, then $s_\lambda(q/u) = 0$.
\end{lemma}

We recall the following.

\begin{proposition}[\!\! {\cite[Theorem 5.1]{SaganSwanson2024}}]\label{prop:OSP-inv}
For $n, k \geq 0$, we have
    \begin{equation}
        [k]_q! S[n,k] = \sum_{w \in \OSP(n,k)}q^{\inv(w)}.
    \end{equation}
\end{proposition}

In other words,
\begin{equation}\label{eq:OSP-equation}
    \langle q^i \rangle [k]_q! S[n,k] = \#\{w \in \OSP(n,k)\, | \, \inv(w) = i \}.
\end{equation}



Now we recall the involution $\phi$ of Sagan--Swanson \cite[Section 5.1]{SaganSwanson2024}.
Given an ordered set partition $w=(S_1/S_2/\ldots /S_k)$, the ordered set partition $w$ is called \textbf{splittable} at $M = \max S_i$, for a fixed $i \in\{1,\ldots,k\}$, if $|S_i|>1$. 
The ordered set partition $w$ is called \textbf{mergeable} at $M= \max S_i$, for a fixed $i \in\{1,\ldots,k-1\}$, if $|S_i|=1$ and $M>\max S_{i+1}$.

Define a function $\phi$ on ordered set partitions of $[n]$.
Given $w=(S_1/S_2/\ldots /S_k)$, find the largest $M$ such that $w$ is either splittable or mergeable at $M$, if such an $M$ exists. We say that
\begin{itemize}
    \item $w$ is type I if $w$ is mergeable at $M$;
    \item $w$ is type II if $w$ is splittable at $M$;
    \item $w$ is type III if no $M$ exists.
\end{itemize}
Then 
\begin{equation}\label{eq:phi-mergeable-splittable}
\phi(w)=\begin{cases}
   (S_1/S_2/\ldots /S_{i-1}/S_i\cup S_{i+1}/S_{i+2}/\ldots  /S_k) & \text{if $w$ is type I}, \\
   (S_1/S_2/\ldots /S_{i-1}/\{M\}/S_i-\{M\}/S_{i+1}/\ldots  /S_k) & \text{if $w$ is type II}, \\
   w & \text{if $w$ is type III}.
\end{cases}
\end{equation}
It follows that $\phi$ is an involution.
As merging removes one inversion and one block
and splitting adds one inversion and one block then if $w$ has $\blocks(w)$ blocks and $\inv(w)$ inversions, then $n-\blocks(w)-\inv(w)$ is preserved by $\phi$.
Finally $\phi$ is a sign-reversing involution on all ordered set partitions (which are not fixed points) with the definition $\sgn(w) = (-1)^{n- \blocks(w)}$.

\subsection{Proofs}
In this subsection, we prove the desired combinatorial formula. First, we express the universal series coefficients $c_{\lambda}(n)$ that we need in terms of a signed enumeration of certain ordered set partitions.

\begin{lemma}\label{lem:c_{(a,1^b)}}
    For $a \geq 1$ and $b \geq 0$, the coefficients $c_{(a,1^b)}(n)$ from equation~(\ref{eq:coeffs-Hilbert}) may be computed by
    \begin{equation}\label{eq:first-formula-coeffs}
        c_{(a,1^b)}(n) = \sum_{i=0}^b (-1)^i \#\{w \in \OSP(n,n-b+i)\, | \, \inv(w) = a+i \}.
    \end{equation}
\end{lemma}

\begin{proof}
By equations~(\ref{eq:coeffs-Hilbert}) and~(\ref{eq:hilbert-rhoades-wilson}), we have that
\begin{equation}
    \sum_{\lambda \in P(1,1,n)} c_{\lambda}(n) s_\lambda(q/u) = \sum_{d=0}^n [d]_q!S[n,d] u^{n-d}.
\end{equation}
Then by Lemma~\ref{lem:schur-1-1}, we have
\begin{equation}
    1 + \sum_{\substack{a \geq 1,\\ b \geq 0}} c_{(a,1^b)}(n)(q^au^b+q^{a-1}u^{b+1}) = \sum_{d=0}^n [d]_q!S[n,d] u^{n-d}.
\end{equation}
By equation~(\ref{eq:OSP-equation}),
\begin{equation}
    1 + \sum_{\substack{a \geq 1,\\ b \geq 0}} c_{(a,1^b)}(n)(q^au^b+q^{a-1}u^{b+1}) = \sum_{d=0}^n  \sum_{i \geq 0} q^i \# \{w \in \OSP(n,d)\, |\, \inv(w) = i \} u^{n-d}.
\end{equation}
By extracting the coefficients of $q^a u^b$ from both sides, we find that, for $b \geq 1$,
\begin{equation}
    c_{(a,1^b)}(n)+c_{(a+1,1^{b-1})}(n) = \#\{w\in \OSP(n,n-b)\ | \ \inv(w)=a\}
\end{equation}
and for $b=0$,
\begin{equation}
    c_{(a)}(n) = \#\{w\in \OSP(n,n)\ | \ \inv(w)=a\}.
\end{equation}
Finally, we solve for $c_{(a,1^b)}(n)$ to establish equation~(\ref{eq:first-formula-coeffs}).


\end{proof}

We desire to improve upon Lemma~\ref{lem:c_{(a,1^b)}} to get a manifestly positive formula for the coefficients $c_{(a,1^b)}(n)$. 

For $0 \leq b \leq {\lfloor \frac{\ell-1}{2}\rfloor}$, define the set
\begin{equation}
S_{n,\ell,b}=\bigcup_{c=0}^b  \{w\in\OSP(n,n-c)\, |\, \inv(w)=\ell-b-c\}.
\end{equation}
This set is a reindexing of the union of the sets that appears in equation~(\ref{eq:first-formula-coeffs}).

While $\phi$ is an involution on the set of all ordered set partitions, it does not necessarily fix $S_{n,\ell,b}$. 
We describe when this happens in the following lemma.
\begin{lemma}
    If $w \in S_{n,\ell,b}$, then:
\begin{enumerate}[(i)]
    \item if $w$ is of type I, $w \in \OSP(n,n-c)$ and $\inv(w) = \ell-2c$, then $\phi(w) \not\in S_{n,\ell,b}$;
    \item if $w$ is otherwise of type I, then $\phi(w) \in S_{n,\ell,b}$;
    \item if $w$ is of type II, then $\phi(w) \in S_{n,\ell,b}$;
    \item $w$ cannot be of type III.
\end{enumerate}
\end{lemma}

\begin{proof}
Let $b$ be fixed. The condition on $\inv(w)$ for $S_{n,\ell,b}$ requires that $\ell - 2b \leq \inv(w) \leq \ell-b$, by considering the values $c$ can take on.
\begin{enumerate}[(i)]
    \item If $w$ is of type I, $w \in \OSP(n,n-c)$ and $\inv(w) = \ell-2c$, this means that $c=b$, so $\inv(w) = \ell-2b$. As $\phi$ applied to an ordered set partition of type I lowers the number of inversions by 1, $\inv(\phi(w)) = \ell-2b-1$, so $\phi(w) \not \in S_{n,\ell,b}$.
    \item If $w$ is otherwise of type I, then $c < b$. We get that $\inv(\phi(w))= \ell - b - c -1 \geq \ell-2b$. Thus $\phi(w) \in S_{n,\ell,b}$.
    \item If $w$ is of type II, then $\blocks(w) < n$, so $c > 0$, which implies that $\inv(w) < \ell-b$. Thus $\inv(\phi(w)) \leq \ell - b$ and $\phi(w) \in S_{n,\ell,b}$.
    \item Suppose towards a contradiction that  $w \in S_{n,\ell,b}$ is of type III, then $c=0$ and $\ell = b$. Since $\ell \geq 1$, then $b > {\lfloor \frac{\ell-1}{2}\rfloor}$, contradicting that $b \leq {\lfloor \frac{\ell-1}{2}\rfloor}$.
\end{enumerate}
\end{proof}


Let us collect those ordered set partitions in $S_{n,\ell,b}$ which are not paired under $\phi$ as 
\begin{equation}
    T_{n,\ell,b}=
\{w\in \OSP(n,n-b) \, |\, \text{$w$ of type I and}\ \inv(w)=\ell-2b\}.
\end{equation}


Now we are ready to state the answer to the problem.

\begin{theorem}\label{thm:improved-coefficient-formula}
        For $a \geq 1$ and $b \geq 0$, the coefficients $c_{(a,1^b)}(n)$ from equation~(\ref{eq:coeffs-Hilbert}) may be computed in a manifestly positive manner by
    \begin{equation}\label{eq:second-formula-coeffs}
        c_{(a,1^b)}(n) = \#\{w \in \OSP(n,n-b)\, | \, \inv(w) = a,\ w \text{ is type I} \}.
    \end{equation}
\end{theorem}

\begin{proof}
    
Use the $\phi$ involution of Sagan--Swanson on 
\begin{equation}
S_{n,a+2b,b} = \bigcup_{i=0}^b \{ w\in \OSP(n,n-b+i)\, | \, \inv(w) = a+i\}.   
\end{equation}
Those points in $S_{n,a+2b,b}$ which are not paired under $\phi$ are
\begin{equation}
    T_{n,a+2b,b} = \{w \in \OSP(n,n-b)\, | \, \inv(w) = a, \ w \text{ is type I} \}.
\end{equation}
\end{proof}

\subsection{Example}
In this subsection, we give an example of our solution at $n=4$.

\begin{figure}[h]
\begin{tabular}{l|l|llllll}
\text{Coefficient} & \text{Value} & \text{OSPs}         \\ \hline
$c_{(1)}(4)$       & 3            & $2|1|3|4$  & $1|3|2|4$ & $1|2|4|3$ &           &           &           \\
$c_{(1,1)}(4)$     & 3            & $1|4|23$   & $12|4|3$  & $3|12|4$  &           &           &           \\
$c_{(1,1,1)}(4)$   & 1            & $4|123$    &           &           &           &           &           \\
$c_{(2)}(4)$       & 5            & $2|1|4|3$  & $1|3|4|2$ & $2|3|1|4$ & $3|1|2|4$ & $1|4|2|3$ &           \\
$c_{(2,1)}(4)$     & 5            & $4|1|23$   & $13|4|2$  & $2|4|13$  & $3|4|12$  & $4|12|3$  &           \\
$c_{(3)}(4)$       & 6            & $2|3|4|1$  & $3|2|1|4$ & $2|4|1|3$ & $3|1|4|2$ & $4|1|2|3$ & $1|4|3|2$ \\
$c_{(3,1)}(4)$     & 4            & $4|13|2$   & $23|4|1$  & $4|2|13$  & $4|3|12$  &           &           \\
$c_{(4)}(4)$       & 5            & $4|1|3|2$  & $4|2|1|3$ & $3|4|1|2$ & $3|2|4|1$ & $2|4|3|1$ &           \\
$c_{(4,1)}(4)$     & 1            & $4|23|1$   &           &           &           &           &           \\
$c_{(5)}(4)$       & 3            & $4|2|3|1$  & $4|3|1|2$ & $3|4|2|1$ &           &           &           \\
$c_{(6)}(4)$       & 1            & $4|3|2|1$  &           &           &           &           &          
\end{tabular}
\caption{A table showing the coefficients $c_\lambda(n)$ and their contributing ordered set partitions (OSPs) given by Theorem~\ref{thm:improved-coefficient-formula}. Aside from $c_\varnothing(4)=1$, for all other $\lambda$ not covered in the table, those coefficients are $c_\lambda(4)=0$.}
\end{figure}

One can compute that
\begin{align*}
    \Hilb(R_n^{(1,1)};q;u) &=  (q^6+3q^5+5q^4+6q^3+5q^2+3q+1)+ (q^5+4q^4+9q^3+11q^2+8q+3)u\\
    &\quad+ (q^3+4q^2+6q+3)u^2 +u^3\nonumber\\
    &= 1+ c_{(1)}(4)s_{(1)}(q/u) + c_{(2)}(4)s_{(2)}(q/u) + c_{(3)}(4)s_{(3)}(q/u) + c_{(4)}(4)s_{(4)}(q/u)\\ 
    &\quad+ c_{(5)}(4)s_{(5)}(q/u) + c_{(6)}(4)s_{(6)}(q/u) + c_{(1,1)}(4)s_{(1,1)}(q/u) + c_{(2,1)}(4)s_{(2,1)}(q/u)\nonumber\\ 
    &\quad + c_{(3,1)}(4)s_{(3,1)}(q/u) + c_{(4,1)}(4)s_{(4,1)}(q/u) + c_{(1,1,1)}(4)s_{(1,1,1)}(q/u).\nonumber
\end{align*}

\begin{thebibliography}{99}
 \bibitem{Bergeron2020}   Fran\c{c}ois Bergeron, The bosonic-fermionic diagonal coinvariant modules conjecture, preprint (2020), \url{https://arxiv.org/abs/2005.00924}.
\bibitem{ChengWangBook} Shun-Jen Cheng and Weiqiang Wang, Dualities and representations of Lie superalgebras, Graduate Studies in
Mathematics, vol. 144, American Mathematical Society, Providence, RI, 2012.
\bibitem{Lentfer-Supersymmetry} John Lentfer, Diagonal supersymmetry for coinvariant rings, preprint (2025), \url{https://arxiv.org/abs/2505.14885}.
\bibitem{RhoadesWilson2023} Brendon Rhoades and Andrew Timothy Wilson, The Hilbert series of the superspace coinvariant ring, Forum Math.
Pi 12 (2024), 35 (English), Id/No e16.
\bibitem{SaganSwanson2024} Bruce E. Sagan and Joshua P. Swanson, $q$-Stirling numbers in type $B$, Eur. J. Comb. 118 (2024), 35, Id/No 103899
\end{thebibliography}


\end{document}